% !TEX root = ../Main.tex
\chapter{拆项法}

数列求和题里 80\% 的关键是\textbf{裂项}——把一个"难以直接求和"的项拆成两个相邻项的差，累加时中间部分消去，只剩首尾两项。按通项结构可分三类：等差型、等比型、等差 + 等比型。此外还有"和项"（每项自带正负号相间）需要\textbf{两两配对再求和}。

\section{裂项：等差型}

\state{等差型裂项的基本模板}{
    若 $a_n = \dfrac{1}{(k_{n+1})(k_n)}$，其中 $k_n = k_1 + (n-1)d$ 为等差数列，则
    \[
    a_n = \frac{1}{d}\!\left(\frac{1}{k_n} - \frac{1}{k_{n+1}}\right).
    \]
    \textbf{这个 $\tfrac{1}{d}$ 因子容易漏！}
}

\ex[等差型裂项]{
    数列 $\{a_n\}$ 满足 $a_n = \dfrac{1}{(2n+1)(2n-1)}$。求前 $n$ 项和 $T_n$。
}

\sol{
    \[
    a_n = \frac{1}{(2n+1)(2n-1)} = \frac{\tfrac{1}{2}\bigl[(2n+1)-(2n-1)\bigr]}{(2n+1)(2n-1)} = \frac{1}{2}\!\left(\frac{1}{2n-1} - \frac{1}{2n+1}\right).
    \]
    累加：
    \begin{align*}
    T_n &= \frac{1}{2}\!\left[\left(1 - \frac{1}{3}\right) + \left(\frac{1}{3} - \frac{1}{5}\right) + \cdots + \left(\frac{1}{2n-1} - \frac{1}{2n+1}\right)\right]\\
    &= \frac{1}{2}\!\left(1 - \frac{1}{2n+1}\right) = \frac{n}{2n+1}.
    \end{align*}
}

\section{裂项：等比型}

\state{等比型裂项的基本模板}{
    若分子含 $q^n$、分母含两项等比差，常可用恒等式
    \[
    q^n = \frac{q^{n+1} - q^n}{q - 1}
    \]
    把分子"拆"到分母里去，裂为相邻两项之差。
}

\ex[等比型裂项 + 主元法]{
    数列 $\{b_n\}$ 通项 $b_n = \dfrac{2^n}{(2^{n+1} - 1)(2^n - 1)}$，前 $n$ 项和为 $S_n$。若 $\forall n \in \mathbb N^+,\ a \in [-1, 1]$，均有 $x^2 - ax - 2 S_n \ge 0$ 成立，求 $x$ 的取值范围。
}

\sol{
    \textbf{第一步：裂 $b_n$。} 利用 $2^n = (2^{n+1} - 1) - (2^n - 1)$：
    \[
    b_n = \frac{(2^{n+1} - 1) - (2^n - 1)}{(2^{n+1} - 1)(2^n - 1)} = \frac{1}{2^n - 1} - \frac{1}{2^{n+1} - 1}.
    \]
    于是
    \[
    S_n = \frac{1}{2^1 - 1} - \frac{1}{2^{n+1} - 1} = 1 - \frac{1}{2^{n+1} - 1}.
    \]

    \textbf{第二步：$n \to \infty$ 的极限思想。} 由 $S_n < 1$ 且 $S_n \to 1$，所以 $\forall n$，$2 S_n < 2$ 且 $2 S_n \to 2$。要使 $x^2 - a x - 2 S_n \ge 0$ 对所有 $n$ 都成立，\textbf{最严格的约束是 $n \to \infty$}，即
    \[
    x^2 - a x - 2 \ge 0 \quad \text{对所有 } a \in [-1, 1] \text{ 恒成立}.
    \]

    \textbf{第三步：主元法——把 $a$ 当主元，$x$ 当参数。} 式子
    \[
    g(a) := -x \cdot a + (x^2 - 2)
    \]
    是 $a$ 的一次函数。要在 $[-1, 1]$ 上恒 $\ge 0$，只需两端点 $\ge 0$：
    \[
    \begin{cases}
    g(-1) = x + x^2 - 2 \ge 0 \iff (x+2)(x-1) \ge 0 \iff x \in (-\infty, -2] \cup [1, +\infty), \\
    g(1) = -x + x^2 - 2 \ge 0 \iff (x-2)(x+1) \ge 0 \iff x \in (-\infty, -1] \cup [2, +\infty).
    \end{cases}
    \]
    取交集：
    \[
    x \in (-\infty, -2] \cup [2, +\infty).
    \]
}

\section{裂项：等差 + 等比型}

\state{等差 + 等比型裂项的原则}{
    对于 $\{a_n\}$ 形如 $a_n = \dfrac{(\text{等差})_n \cdot q^n}{(\text{等差})_n \cdot (\text{等差})_{n+1}}$ 这种"等差与等比混合"的通项，裂项后的形式应是
    \[
    a_n = c_{n+1} - c_n,\qquad c_n = \frac{q^n}{\text{某等差}}.
    \]
    裂项后 $\{c_n\}$ 的\textbf{分子应仅差 $q^n$}（因为 $q^n$ 的比是 $q$，天然适合两相邻项之差），而\textbf{除 $q^n$ 外只有与 $n$ 无关的因子}。
}

\ex[等差 + 等比型]{
    数列 $\{a_n\}$ 满足 $a_n = \dfrac{(2n-3)\cdot 2^n}{4 n^2 - 1}$。若 $S_n$ 为 $\{a_n\}$ 的前 $n$ 项和，求 $S_n$。
}

\sol{
    \textbf{分母凑分子：} $4n^2 - 1 = (2n+1)(2n-1)$。关键恒等式
    \[
    2n - 3 = 2(2n-1) - (2n+1).
    \]
    代入：
    \[
    a_n = \frac{\bigl[2(2n-1) - (2n+1)\bigr] \cdot 2^n}{(2n+1)(2n-1)} = \frac{2^{n+1}}{2n+1} - \frac{2^n}{2n-1}.
    \]

    设 $c_n = \dfrac{2^n}{2n-1}$，则 $c_{n+1} = \dfrac{2^{n+1}}{2n+1}$，且 $a_n = c_{n+1} - c_n$。故
    \[
    S_n = c_{n+1} - c_1 = \frac{2^{n+1}}{2n+1} - 2.
    \]
}

\section{和项：正负号相间}

\state{两两配对}{
    若通项含 $(-1)^n$，直接累加会忽上忽下，难以看出规律。\textbf{诀窍是把奇、偶两项配对求和}，让振荡"局部相消"。分类讨论 $n$ 的奇偶即可。
}

\ex[带 $(-1)^n$ 的和项]{
    数列 $\{c_n\}$ 通项 $c_n = (-1)^n \dfrac{4n}{(2n+1)(2n-1)}$，$\{c_n\}$ 的前 $n$ 项和为 $T_n$。求 $T_n$。
}

\sol{
    \textbf{先裂项：} 由 $4n = (2n+1) + (2n-1)$，
    \[
    c_n = (-1)^n \!\left(\frac{1}{2n-1} + \frac{1}{2n+1}\right).
    \]

    \textbf{奇偶配对：} 计算 $c_n + c_{n+1}$（其中 $n$ 为奇数，$n+1$ 为偶数）：
    \[
    c_n + c_{n+1} = -\!\left(\frac{1}{2n-1} + \frac{1}{2n+1}\right) + \left(\frac{1}{2n+1} + \frac{1}{2n+3}\right) = -\frac{1}{2n-1} + \frac{1}{2n+3}.
    \]
    这是个漂亮的"跨两步"裂项结构。

    \textbf{情形一：$n$ 为奇数（$n = 2k-1$）。} 前 $n$ 项配成 $\tfrac{n-1}{2}$ 对加一个剩余项 $c_n$：
    \[
    T_n = \underbrace{(c_1 + c_2) + (c_3 + c_4) + \cdots + (c_{n-2} + c_{n-1})}_{\text{每对都是裂项差}} + c_n.
    \]
    对子求和（裂项累加）：
    \[
    \sum_\text{pairs} = \left(-\frac{1}{1} + \frac{1}{5}\right) + \left(-\frac{1}{5} + \frac{1}{9}\right) + \cdots + \left(-\frac{1}{2n-5} + \frac{1}{2n-1}\right) = -1 + \frac{1}{2n-1}.
    \]
    加上 $c_n = -\dfrac{1}{2n-1} - \dfrac{1}{2n+1}$（因 $n$ 为奇数，$(-1)^n = -1$）：
    \[
    T_n = -1 + \frac{1}{2n-1} - \frac{1}{2n-1} - \frac{1}{2n+1} = -1 - \frac{1}{2n+1} = -\frac{2n+2}{2n+1}.
    \]

    \textbf{情形二：$n$ 为偶数。} 全部配对：
    \[
    T_n = (c_1 + c_2) + \cdots + (c_{n-1} + c_n) = -1 + \frac{1}{2n+1} = -\frac{2n}{2n+1}.
    \]

    综合：
    \[
    T_n = \begin{cases} -\dfrac{2n+2}{2n+1}, & n\text{ 为奇数}, \\[4pt] -\dfrac{2n}{2n+1}, & n\text{ 为偶数}. \end{cases}
    \]
}

\rmkb{
    本章四类裂项背后是\textbf{同一个朴素想法}："让累加后中间消，只留首尾"。具体形式因 $a_n$ 的"骨架"而异，但识别它们只需做一件事——\textbf{看分母/分子的因子能不能用相邻两项的差表达}。这也是为什么\textbf{等差 + 等比型}要先"分母凑分子"：让分子变成两个相邻 $(2n\pm 1)$ 的线性组合。
}
